% SPDX-FileCopyrightText: © 2026 Idiap Research Institute <contact@idiap.ch>
%
% SPDX-FileContributor: Andre Freitas <andre.freitas@idiap.ch>
% Neuro-symbolic AI Group
%
% SPDX-FileContributor: Danilo Gusicuma <danilo.gusicuma@idiap.ch>
% Neuro-symbolic AI Group
%
% SPDX-License-Identifier: MIT
\subsection{Following one premise through four Prism steps}
\label{sec:step-walkthrough}
To make the distinction between an interface and a reasoning procedure concrete,
consider a constructed task: establish diffusion decay with one positive rate for
every graph in a specified family of finite connected undirected graphs. Let $L$
be the symmetric combinatorial Laplacian, and restrict $u(0)$ to the mean-zero
subspace. The desired conclusion is
\[
 \|u(t)\|_2\leq e^{-ct}\|u(0)\|_2,
 \qquad \dot u=-Lu,\quad t\geq0,
\]
with the same $c>0$ throughout the family. The four native occurrences of the
abductive work loop can now be read as four distinct questions. Their internal
procedures below are source-guided prescriptions and the values are illustrative;
this is not a recorded execution.

\begin{enumerate}[leftmargin=*]
\item \textbf{Formation: what premise is missing?}
The input is the unclosed uniform-decay obligation. The linked abductive procedure
first attempts deduction to locate the missing support, then proposes small sets
of additional premises and compares their coverage, compatibility and falsifiers.
Here the useful candidate is a common lower bound $\lambda_1(L)\geq c$ on the
first positive eigenvalue. The output is this \emph{proposed} premise, with its
family-wide quantifier. It changes the proof question from ``derive decay somehow''
to ``does this family have a common gap?'' It does not assert that the gap exists.

\item \textbf{Derivation: what follows if that premise holds?}
The next occurrence consumes the proposed route. Expanding the initial vector in
orthonormal positive-eigenvalue modes gives
\[
 u(t)=\sum_{j\geq1}\alpha_j e^{-\lambda_jt}v_j,
 \qquad
 \|u(t)\|_2^2=\sum_{j\geq1}|\alpha_j|^2e^{-2\lambda_jt}
 \leq e^{-2ct}\|u(0)\|_2^2.
\]
The output is the conditional decay claim. Connectedness and mean-zero data remove
the zero mode; symmetry supplies the orthonormal spectral decomposition. The
assumption relation licenses using the proposed common gap as a premise, while
leaving its proof burden intact. Dropping that condition would change the claim.

\item \textbf{Refutation: which commitment can fail?}
The test receives the derived claim, including its stated premises. If the family
is all connected graphs, cycle graphs provide a discriminating boundary case:
\[
 \lambda_1(C_n)=2-2\cos(2\pi/n)\longrightarrow0.
\]
Thus connectedness does not supply the proposed positive common bound. The finding
attacks the route's attempted uniform strengthening, not the conditional spectral
calculation. If the input were restricted to a family with an independently proved
common gap, this particular objection would no longer defeat its premise.

\item \textbf{Calibration: what should change next?}
The final occurrence consumes both the retained derived claim and its refutation
assessment through named fields. It can therefore preserve the valid conditional
estimate while recording why the unrestricted route fails. A useful continuation
is to seek a stronger structural assumption or restrict the family; repeating the
same connectedness argument adds no support. The original workflow policy provides
reframe, retarget, release and stop routes. The native sequence returns a routing
assessment and does not execute those downstream branches.
\end{enumerate}

This reconstruction makes the informational change at each boundary visible:
an obligation becomes a candidate premise, the premise becomes a conditional
consequence, a boundary case becomes a targeted criticism, and the retained
claim--criticism pair determines the next question. Appendix~\ref{app:granular-prism-workflows}
provides this step-level account for every declared occurrence in all nineteen
workflows, including unresolved producers, parallel inputs and handoff gaps.
